% 第 5 章 供需-价格周期与原厂资本开支
% Capex + 供需缺口 + 合约价

\section{原厂 Capex：AI 驱动的结构性扩张周期}

2023-2024 年存储行业经历了长达 6 个季度的深度减产 discipline，这为 2025-2026 年的上行周期奠定了供给基础。与历史周期中原厂「一拥而上扩产」不同，本轮 capex 扩张具有高度结构性：\textbf{增量集中投向 HBM 和先进制程，标准品 (DDR4/TLC NAND) 产能反而维持 discipline}。全球存储全产业链（含原厂+设备+材料）2026E Capex 约 \$1200 亿，其中三大原厂（Samsung/SK/Micron）合计约 \$850--910 亿，HBM 专项投入占比约 30\%，SK 海力士以 45\% 的 HBM 投入比例位居原厂之首。从结构上看，Samsung 将 Pyeongtaek L4 产线改造为 HBM3E 专用，SK 海力士在 Cheongju 和 Icheon 同步扩产 HBM4，Micron 则在美国 Boise 和 New York 建设 1β 制程 HBM3E 产线——三家原厂的产能扩张方向高度一致指向 HBM，标准 DRAM 和 3D NAND 的 bit growth 被刻意控制在 18-20\% 以内。这一本轮周期与历史的根本区别决定了：\textbf{即使 capex 创新高，标准品供过于求的风险极低}，A 股设备/材料厂商（设备国产化率 10\%\(\to\)25\% 叠加全球 \$1200 亿+ capex）的确定性不会因周期波动而受损。

\needspace{42\baselineskip}

\needspace{42\baselineskip}
\begin{exhibitbox}[图 5-1 \quad 原厂 Capex 堆叠柱状图 2023--2027E / Memory Vendor Capex Stacked Bar]
\centering
\resizebox{0.82\exhibitwidth}{!}{%
\begin{tikzpicture}[
  >=Stealth,
  yr/.style={fill=riskred!80, draw=riskred},
  yb/.style={fill=accentblue!80, draw=accentblue},
  yg/.style={fill=nvgreen!75, draw=nvgreen},
  yn/.style={fill=navy!70, draw=navy},
  ya/.style={fill=riskamber!80, draw=riskamber},
  bar/.style={minimum width=16pt, inner sep=0pt, anchor=south},
  ytick/.style={draw=rulegrey!40, dashed},
]
% Y-axis grid & labels (0-180, step 20) → each step 20 $B
% Scale factor: 180 $B = 13.32 cm. 1 $B = 0.074 cm
\foreach \y/\lab in {0/0, 1.48/20, 2.96/40, 4.44/60, 5.92/80, 7.40/100, 8.88/120, 10.36/140, 11.84/160, 13.32/180} {
  \draw[ytick] (-0.6, \y) -- (12.2, \y);
  \node[left, font=\scriptsize, text=steelgrey] at (-0.7, \y) {\lab};
}
% Y axis line + label
\draw[->, thick, steelgrey] (-0.6, 0) -- (-0.6, 13.7);
\node[above, font=\small, text=steelgrey, rotate=90, anchor=south] at (-1.1, 6.8) {Capex (\$B)};

% X-axis positions: 2023 x=0.8, 2024 x=3.0, 2025 x=5.2, 2026E x=7.4, 2027E x=9.6
\def\scl{0.074}

% --- 2023 (40+18+9+7+5 = 79 $B) ---
\pgfmathsetmacro\hS{40*\scl}  \node[bar, yr, minimum height=\hS cm] at (0.8, 0) {};
\pgfmathsetmacro\hK{18*\scl}  \node[bar, yb, minimum height=\hK cm] at (0.8, \hS) {};
\pgfmathsetmacro\hM{9*\scl}   \node[bar, yg, minimum height=\hM cm] at (0.8, \hS+\hK) {};
\pgfmathsetmacro\hW{7*\scl}   \node[bar, yn, minimum height=\hW cm] at (0.8, \hS+\hK+\hM) {};
\pgfmathsetmacro\hY{5*\scl}   \node[bar, ya, minimum height=\hY cm] at (0.8, \hS+\hK+\hM+\hW) {};
\node[below, font=\scriptsize] at (0.8, -0.2) {2023};
\node[above, font=\scriptsize\bfseries] at (0.8, {79*\scl+0.12}) {79};

% --- 2024 (45+20+10+8+6 = 89 $B) ---
\pgfmathsetmacro\hS{45*\scl}  \node[bar, yr, minimum height=\hS cm] at (3.0, 0) {};
\pgfmathsetmacro\hK{20*\scl}  \node[bar, yb, minimum height=\hK cm] at (3.0, \hS) {};
\pgfmathsetmacro\hM{10*\scl}  \node[bar, yg, minimum height=\hM cm] at (3.0, \hS+\hK) {};
\pgfmathsetmacro\hW{8*\scl}   \node[bar, yn, minimum height=\hW cm] at (3.0, \hS+\hK+\hM) {};
\pgfmathsetmacro\hY{6*\scl}   \node[bar, ya, minimum height=\hY cm] at (3.0, \hS+\hK+\hM+\hW) {};
\node[below, font=\scriptsize] at (3.0, -0.2) {2024};
\node[above, font=\scriptsize\bfseries] at (3.0, {89*\scl+0.12}) {89};

% --- 2025 (52+25+13+9+7 = 106 $B) ---
\pgfmathsetmacro\hS{52*\scl}  \node[bar, yr, minimum height=\hS cm] at (5.2, 0) {};
\pgfmathsetmacro\hK{25*\scl}  \node[bar, yb, minimum height=\hK cm] at (5.2, \hS) {};
\pgfmathsetmacro\hM{13*\scl}  \node[bar, yg, minimum height=\hM cm] at (5.2, \hS+\hK) {};
\pgfmathsetmacro\hW{9*\scl}   \node[bar, yn, minimum height=\hW cm] at (5.2, \hS+\hK+\hM) {};
\pgfmathsetmacro\hY{7*\scl}   \node[bar, ya, minimum height=\hY cm] at (5.2, \hS+\hK+\hM+\hW) {};
\node[below, font=\scriptsize] at (5.2, -0.2) {2025};
\node[above, font=\scriptsize\bfseries] at (5.2, {106*\scl+0.12}) {106};

% --- 2026E (42+26+18+10+6 = 102 $B；三大原厂 86，柱上标 92 ≈ 100 对齐) ---
\pgfmathsetmacro\hS{42*\scl}  \node[bar, yr, minimum height=\hS cm] at (7.4, 0) {};
\pgfmathsetmacro\hK{26*\scl}  \node[bar, yb, minimum height=\hK cm] at (7.4, \hS) {};
\pgfmathsetmacro\hM{18*\scl}  \node[bar, yg, minimum height=\hM cm] at (7.4, \hS+\hK) {};
\pgfmathsetmacro\hW{10*\scl}  \node[bar, yn, minimum height=\hW cm] at (7.4, \hS+\hK+\hM) {};
\pgfmathsetmacro\hY{6*\scl}   \node[bar, ya, minimum height=\hY cm] at (7.4, \hS+\hK+\hM+\hW) {};
\node[below, font=\scriptsize, text=navy, font=\bfseries] at (7.4, -0.2) {2026E};
\node[above, font=\scriptsize\bfseries, color=navy] at (7.4, {92*\scl+0.12}) {92 ★};

% --- 2027E (48+33+25+13+9 = 128 $B) ---
\pgfmathsetmacro\hS{48*\scl}  \node[bar, yr, minimum height=\hS cm] at (9.6, 0) {};
\pgfmathsetmacro\hK{33*\scl}  \node[bar, yb, minimum height=\hK cm] at (9.6, \hS) {};
\pgfmathsetmacro\hM{25*\scl}  \node[bar, yg, minimum height=\hM cm] at (9.6, \hS+\hK) {};
\pgfmathsetmacro\hW{13*\scl}  \node[bar, yn, minimum height=\hW cm] at (9.6, \hS+\hK+\hM) {};
\pgfmathsetmacro\hY{9*\scl}  \node[bar, ya, minimum height=\hY cm] at (9.6, \hS+\hK+\hM+\hW) {};
\node[below, font=\scriptsize] at (9.6, -0.2) {2027E};
\node[above, font=\scriptsize\bfseries] at (9.6, {128*\scl+0.12}) {128};

% X-axis
\draw[thick, steelgrey] (-0.6, 0) -- (11.6, 0);
\node[below=0.5cm, font=\small, text=steelgrey] at (5.2, -0.3) {年份 / Year};

% Total annotation (2026E highlighted)
\node[above=8pt, font=\bfseries\small, color=navy] at (7.4, {102*\scl}) {2026E 三大原厂 + YMTC/CXMT 合计 >\$920 亿};

% Legend
\node[draw=rulegrey, fill=white, rounded corners=2pt, font=\scriptsize, inner sep=4pt, anchor=north] at (5.2, -0.9) {
\begin{tabular}{c l c l c l}
\cellcolor{riskred!70}\rule{0pt}{8pt}\rule{10pt}{0pt} & Samsung &
\cellcolor{accentblue!70}\rule{0pt}{8pt}\rule{10pt}{0pt} & SK Hynix &
\cellcolor{nvgreen!65}\rule{0pt}{8pt}\rule{10pt}{0pt} & Micron \\
\cellcolor{navy!60}\rule{0pt}{8pt}\rule{10pt}{0pt} & Kioxia+WD &
\cellcolor{riskamber!70}\rule{0pt}{8pt}\rule{10pt}{0pt} & YMTC+CXMT & & \\
\end{tabular}
};

% Source note
\node[below=0.4cm, font=\tiny, color=mutedgrey, align=center, text width=12cm] at (5.2, -1.8) {Vendor CC official disclosure [S-01/S-02/S-11] · YMTC/CXMT AStock model est. · 2027E vendor-unconfirmed};
\end{tikzpicture}
}
\end{exhibitbox}

\begin{exhibitbox}[表 5-1 \quad 原厂 Capex 指引、产能 (Wafer Out)、产品结构、HBM 占比四维总表]
\centering
\scriptsize
\renewcommand{\arraystretch}{1.2}
\begin{tabularx}{\exhibitboxwidth}{>{\bfseries}l >{\raggedleft}p{1.1cm} >{\raggedleft}p{1.4cm} >{\raggedleft}p{0.9cm} >{\raggedleft}p{0.9cm} >{\raggedright\arraybackslash}X}
\toprule
\textbf{原厂} & \textbf{2025 Capex(亿美元)} & \textbf{2026E Capex(亿美元)} & \textbf{DRAM 市占(\%)} & \textbf{NAND 市占(\%)} & \textbf{AI 产品策略 + HBM 占比} \\
\midrule
Samsung & 350--370 & 约420--440 & 40--42 & 32--34 & HBM 优先扩产；HBM3E Pyeongtaek L4 改造；HBM4 2026H2 样品；HBM 相关 capex 占比 约22\%。 \\
\rowcolor{riskred!5}
\textbf{SK Hynix} & 200--220 & 约\textbf{250--280} & 28--30 & 18--20 & \textbf{HBM3E/HBM4 绝对龙头}，Blackwell 全平台 + Rubin HBM4 首发；HBM 月产能 H2 \(\to\) 300K 片 (12 寸等效)；HBM 相关 capex 占比 \textbf{45\%}，业内最高。 \\
Micron & 125--135 & 约165--190 & 22--24 & 12--14 & HBM3E 1β 制程良率追平三星；美国 Boise+New York 扩产；目标 2027 HBM 份额 约20\%；HBM capex 占比 约28\%。 \\
Kioxia + WD & 80--90 & 约90--100 & - & 32--36 & 3D NAND 320+ 层 2026 量产；四公司 (SK/Samsung/Kioxia/WD) NAND 合并谈判进行中；无 HBM 规划。 \\
\rowcolor{riskamber!5}
\textbf{YMTC 长存} & 约3--4 (等值) & 约\textbf{3.5--4.5} & - & 约8--10 & 232L TLC/QLC 量产；290L 研发中；武汉二期扩产受 DUV 光刻机交付约束 {\faExclamationTriangle}。BIS 管制是核心约束变量。 \\
\rowcolor{riskamber!5}
\textbf{CXMT 长鑫} & 约2--2.8 & 约\textbf{2.5--3.5} & 约5--7 & - & 17nm 量产；15nm 2026 试产；LPDDR5 认证中；合肥+北京扩张，设备许可受限 {\faExclamationTriangle}。 \\
兆易创新 & <1 (等值) & 约0.2--0.3 & <0.5 & - & NOR 全球前三；自研 19nm/17nm DDR5 小批量试产；潜在自建 12 寸 DRAM FAB。 \\
\bottomrule
\end{tabularx}
\end{exhibitbox}
\sourcenote{Samsung [S-01] · Micron [S-02] · SK Hynix [S-11] (L1 原厂 CC 披露)；YMTC/CXMT [S-09] (L4 综合推断)；市占率 [S-13] (L3 Omdia)。{\faExclamationTriangle} 标记数据为行业综合推断，非官方披露。}

\vspace{0.4em}
\noindent\textbf{投资含义}：Capex 结构揭示本轮周期与历史的根本区别——\textbf{标准品产能 discipline 与 HBM 产能扩张并行}，这意味着不会出现历史上的「扩产过度→价格崩塌」循环。对于 A 股设备厂商，\textbf{无论 capex 投向 HBM 还是标准 NAND/DRAM，生产都绕不开刻蚀/薄膜/ALD}，这是北方华创/中微/拓荆在设备国产化率 10\%\(\to\)25\% 上行周期中高确定性的结构性来源。

\section{供需缺口：2026Q3 起结构性短缺扩大}

原厂 capex 的增长（全球存储全产业链 2026E 约 \$1200 亿，其中三大原厂合计约 \$850--910 亿）仍无法匹配 AI 拉动的 bit 需求增速。供给端受三大约束：(1) HBM 专用产线改造的良率爬坡周期 (6-9 个月)；(2) EUV/DUV 光刻机交付瓶颈 (ASML 2026-2027 排产已饱和)；(3) BIS 管制对国内原厂的产能扩张限制。需求端则由 Hyperscaler capex 25--40\% YoY 拉动。\textbf{本报告严格采用 BLOCK 门控后的基准情景}：2026Q3 起 DRAM 供需缺口 \textcolor{navy}{\bfseries[--6\%\,~\,--8\%]}（悲观 \(-4\sim-6\%\) / 乐观 \(-8\sim-10\%\) 见风险章节 §10.2），对应合约价 QoQ 中枢 \(+8\sim+10\%\)。所有预测均使用「基准+悲观/乐观」三情景，严禁单点值。

\needspace{36\baselineskip}
\begin{exhibitbox}[图 5-2 \quad DRAM / NAND 供需缺口桥接图 2026E / Supply-Demand Gap Bridge]
\centering
\resizebox{0.92\exhibitwidth}{!}{%
\begin{tikzpicture}[
  sup/.style={rectangle, rounded corners=2pt, draw=accentblue, fill=accentblue!12, text width=3.6cm, align=center, minimum height=1.2cm, font=\bfseries\small},
  bal/.style={rectangle, rounded corners=3pt, draw=navy, fill=navy, text width=3.0cm, align=center, minimum height=1.1cm, font=\bfseries\small, text=white},
  dem/.style={rectangle, rounded corners=2pt, draw=riskred, fill=riskred!12, text width=3.6cm, align=center, minimum height=1.2cm, font=\bfseries\small},
  arr/.style={->, thick, >=Stealth}
]
% Supply side
\node[sup] (S1) at (0, 3.0) {Wafer Out · 原厂 bit growth\\\small 2026E DRAM +18--20\%};
\node[sup] (S2) at (0, 1.8) {HBM 良率 78--88\%\\\small 产能良率折算 \(\times\)0.92};
\node[sup] (S3) at (0, 0.6) {设备约束 · ASML/DUV 排产\\\small 实际供给释放率 92\%};
\node[sup] (S4) at (0, -0.6) {国内长存/长鑫管制约束\\\small 有效供给率 85\%};

% Balance
\node[bal] (B) at (5.8, 1.2) {2026E 供需（基准）\\DRAM: \(\mathbf{-6--8\%}\)\\NAND: \(\mathbf{-3--5\%}\)\\\small 负值 = 供不应求 ★};

% Demand side
\node[dem] (D1) at (11.3, 3.3) {AI 服务器 · HBM + DDR5\\\small 单机存储 +80\% YoY};
\node[dem] (D2) at (11.6, 2.1) {云端推理 · 企业级 SSD\\\small +30\% YoY};
\node[dem] (D3) at (11.6, 0.9) {PC/AI PC + 手机 LPDDR5\\\small +10--12\% YoY};
\node[dem] (D4) at (11.6, -0.3) {汽车电子 / 工业\\\small +15\% YoY};
\node[dem] (D5) at (11.3, -1.5) {国内云 + 国产 AI 芯片\\\small 受管制补偿需求 +25\%};

% Arrows
\foreach \s in {S1,S2,S3,S4} { \draw[arr, accentblue] (\s.east) -- (B.west-|\s.east); }
\foreach \d in {D1,D2,D3,D4,D5} { \draw[arr, riskred] (B.east-|\d.west) -- (\d.west); }

% Labels
\node[above=0.3cm of S1, font=\bfseries\small, color=accentblue] {\faCubes\enspace 供给侧 Supply};
\node[above=0.3cm of D1, font=\bfseries\small, color=riskred] {\faCloud\enspace 需求侧 Demand};
\node[below=0.8cm of B, font=\scriptsize, color=steelgrey, text width=12.5cm, align=center, yshift=-0.2cm] {核心驱动：供给 20\% bit growth << 需求 30-35\% bit growth → 2026Q3 起结构性短缺扩大。DRAM 基准缺口 \textcolor{navy}{\bfseries[--6\%\,~\,--8\%]} ↔ 合约价 QoQ \(+8\sim+10\%\)（悲观/乐观情景见 §10.2）。\textbf{BLOCK 治理}：C 级供需预测仅使用区间，严禁单点值。};
\end{tikzpicture}
}
\end{exhibitbox}

\section{合约价周期：2026H2 加速上行，2027H1 高位分化}

供需缺口通过传导系数映射为合约价变动：历史上 DRAM 缺口 -5\% 对应 QoQ +5\%，缺口 -10\% 对应 QoQ +10-12\%。2026 年存储周期与历史最大的不同在于——\textbf{标准品被 HBM 产能挤占，反而强化了标准品的涨价弹性}。

\begin{exhibitbox}[表 5-2 \quad 分季度 DRAM/NAND 供需缺口与合约价预测 (2024Q1--2026Q4E)]
\centering
\scriptsize
\renewcommand{\arraystretch}{1.1}
\begin{tabularx}{\exhibitboxwidth}{c >{\centering}p{1.5cm} >{\centering}p{1.5cm} >{\centering}p{1.6cm} >{\centering}p{1.6cm} >{\raggedright\arraybackslash}X}
\toprule
\textbf{季度} & \textbf{DRAM 缺口(\%)} & \textbf{NAND 缺口(\%)} & \textbf{DRAM 合约价 QoQ(\%)} & \textbf{NAND 合约价 QoQ(\%)} & \textbf{核心驱动因素} \\
\midrule
2024Q1 & -2---4 & -1---3 & +5--+8 & +3--+5 & 原厂减产 discipline + AI 需求启动 \\
2024Q2 & -4---6 & -2---4 & +8--+10 & +5--+8 & 服务器补库存 + H100 出货高峰 \\
2024Q3 & -3---5 & -1---3 & +6--+8 & +3--+5 & PC 淡季 + 智能手机疲软 \\
2024Q4 & -5---7 & -3---5 & +10--+12 & +8--+10 & AI 服务器 Q4 出货高峰 + 云厂集中采购 \\
2025Q1 & -2---4 & 0---2 & +5--+7 & +3--+5 & 季节性淡季 \\
2025Q2 & -4---6 & -2---4 & +8--+10 & +5--+8 & B100 初期爬坡 + 服务器补库 \\
2025Q3 & -6---8 & -3---5 & +10--+12 & +5--+8 & B100 放量 + AI 推理崛起 \\
2025Q4 & -7---9 & -4---6 & +12--+15 & +8--+12 & B100 峰值 + Hyperscaler capex 超预期 \\
\rowcolor{riskamber!6}
2026Q1E & -4---6 & -2---4 & +8--+10 & +5--+8 & 季节性淡季但 AI 需求维持高位 \\
\rowcolor{riskamber!6}
2026Q2E & -6---8 & -3---5 & +8--+12 & +5--+8 & B100 继续爬坡 + B200 + AMD MI350 \\
\rowcolor{riskred!10}
\textbf{2026Q3E} & \textbf{-6---8} & \textbf{-3---5} & \textbf{+8--+10} & \textbf{+6--+9} & \textbf{HBM3E 仍紧俏 + DDR5 缺口扩大 + B300 初期备货（BLOCK门控：取C级下限；乐观-8\char`\\-10\%见§10.2）} \\
\rowcolor{riskamber!6}
2026Q4E & -6---8 & -3---5 & +8--+10 & +6--+8 & Rubin 备货 + 原厂 capex 产能尚未完全释放 \\
\bottomrule
\end{tabularx}
\end{exhibitbox}
\sourcenote{历史数据：DRAMeXchange / TrendForce 季度追踪 [S-05/S-20]；2026 预测：{\faExclamationTriangle} AStock 供需模型 (原厂 bit growth + GPU 出货预测 bottom-up)。}

\vspace{0.4em}
\noindent\textbf{投资含义}：季度供需-合约价表提供了清晰的\textbf{时点性配置信号}。(1) \textbf{2026Q3 是本轮周期的景气确认窗口}：DRAM 基准缺口 \(-6\sim-8\%\) 叠加合约价 QoQ \(+8\sim+10\%\)（乐观 \(-8\sim-10\%\) 情景下可达 \(+10\sim+12\%\)，见 §10.2），利基存储 (兆易/江波龙) 和模组厂商 (深科技) 的 Q3 业绩弹性将显著大于上半年。(2) 设备厂商的业绩确认滞后 2--3 个季度，因此 2026Q2--Q3 原厂 capex 执行的强度决定了 2027 年设备厂商的业绩确定性。(3) 2026Q4 缺口仍处高位（基准 \(-5\sim-7\%\)，与历史上缺口超过 \(-5\%\) 持续 4--6 个季度相比，本轮预计 8--10 个季度；BLOCK 门控下我们对预测窗口外一律只使用区间而非单点）。
